\section{The Lindstr\"{o}m--Gessel--Viennot lemma}

\subsection{Definitions}

\begin{definition}
\label{def.lgv.lattice}
\lean{LGV1.LatticePoint, LGV1.LatticeStep, LGV1.IntegerLatticeAdj, LGV1.integerLattice}
\leantarget
\leanok
We consider the infinite simple digraph with vertex set $\mathbb{Z}^{2}$
(so the vertices are pairs of integers) and arcs
\[
  \left(  i,j\right)  \rightarrow\left(  i+1,j\right)
  \ \ \ \ \ \ \ \ \ \ \text{for all }\left(  i,j\right)  \in\mathbb{Z}^{2}
\]
and
\[
  \left(  i,j\right)  \rightarrow\left(  i,j+1\right)
  \ \ \ \ \ \ \ \ \ \ \text{for all }\left(  i,j\right)  \in\mathbb{Z}^{2}.
\]
The arcs of the first form are called \emph{east-steps} or \emph{right-steps};
the arcs of the second form are called \emph{north-steps} or \emph{up-steps}.

The vertices of this digraph will be called \emph{lattice points} or
\emph{grid points} or simply \emph{points}.

The entire digraph will be denoted by $\mathbb{Z}^{2}$ and called the
\emph{integer lattice} or \emph{integer grid}.

Any path is uniquely determined by its starting point and its step sequence.
\end{definition}

\subsection{Helpers for lattice path infrastructure}

\begin{lemma}
\label{lem.lgv.endpoint-append}
\lean{LGV1.LatticePath.endpoint_append}
\leanhelper
\leanok
Let $p_1, p_2$ be two lattice paths and $A$ a starting point. Then the endpoint of
the concatenation $p_1 \cdot p_2$ starting from $A$ equals the endpoint of $p_2$
starting from the endpoint of $p_1$:
\[
  \operatorname{endpoint}(p_1 \cdot p_2, A) = \operatorname{endpoint}(p_2, \operatorname{endpoint}(p_1, A)).
\]
\end{lemma}

\begin{proof}
By induction on $p_1$: applying the steps of $p_1 \cdot p_2$ starting from $A$
is the same as first applying the steps of $p_1$ (reaching the endpoint of $p_1$)
and then applying the steps of $p_2$.
\end{proof}

\begin{lemma}
\label{lem.lgv.path-take-drop}
\lean{LGV1.LatticePath.isPathFromTo_take_drop}
\leanhelper
\leanok
If a path $p$ goes from $A$ to $B$, then for any $n$, the first $n$ steps of $p$ form a path from $A$
to the intermediate vertex $M$, and
the remaining steps form a path from $M$ to $B$.
\end{lemma}

\begin{proof}
\leanok
Follows from the fact that $p$ equals its first $n$ steps concatenated with its remaining steps,
together with Lemma~\ref{lem.lgv.endpoint-append}.
\end{proof}

\begin{lemma}
\label{lem.lgv.path-append}
\lean{LGV1.LatticePath.isPathFromTo_append}
\leanhelper
\leanok
If path $h$ goes from $A$ to $v$ and path $t$ goes from $v$ to $B$,
then $h \cdot t$ goes from $A$ to $B$.
\end{lemma}

\begin{proof}
\leanok
By Lemma~\ref{lem.lgv.endpoint-append}: $\operatorname{endpoint}(h \cdot t, A) =
\operatorname{endpoint}(t, \operatorname{endpoint}(h, A)) = \operatorname{endpoint}(t, v) = B$.
\end{proof}

\begin{lemma}
\label{lem.lgv.coordsum-endpoint}
\lean{LGV1.LatticePath.coordSum_endpoint}
\leanhelper
\leanok
For any lattice path $p$ starting at $A$, the coordinate sum of the endpoint satisfies
\[
  \operatorname{cs}(\operatorname{endpoint}(p, A)) = \operatorname{cs}(A) + \ell(p),
\]
where $\ell(p)$ is the number of steps in $p$.
\end{lemma}

\begin{proof}
By induction on $p$: each step increases the coordinate sum by exactly $1$.
\end{proof}

\begin{lemma}
\label{lem.lgv.endpoint-x}
\lean{LGV1.endpoint_x_eq}
\leanhelper
\leanok
The x-coordinate of the endpoint of a path $p$ starting at $A$ satisfies
\[
  \operatorname{x}(\operatorname{endpoint}(p, A)) = \operatorname{x}(A) + (\text{\# of east-steps in } p).
\]
\end{lemma}

\begin{proof}
\leanok
By induction: each east-step increments the x-coordinate by $1$, and each north-step
leaves it unchanged.
\end{proof}

\begin{lemma}
\label{lem.lgv.paths-finite}
\lean{LGV1.pathsFromTo_finite}
\leanhelper
\leanok
For any two lattice points $A, B$, the set of lattice paths from $A$ to $B$ is finite.
\end{lemma}

\begin{proof}
\leanok
If $B$ is not componentwise $\ge A$, the set is empty.
Otherwise, every path from $A$ to $B$ has a bounded number of steps
(exactly $\operatorname{cs}(B) - \operatorname{cs}(A)$),
and there are finitely many lists of bounded length over the two-element set
$\{\text{east}, \text{north}\}$.
\end{proof}

\begin{lemma}
\label{lem.lgv.path-coords-nondecreasing}
\lean{LGV1.path_coords_nondecreasing}
\leanhelper
\leanok
If a path $p$ goes from $A$ to $B$, then $\operatorname{x}(A) \le \operatorname{x}(B)$
and $\operatorname{y}(A) \le \operatorname{y}(B)$.
\end{lemma}

\begin{proof}
\leanok
By induction on the path: each east-step increments the x-coordinate and each
north-step increments the y-coordinate.
\end{proof}

\subsection{Counting paths from $(a,b)$ to $(c,d)$}

\begin{lemma}
\label{lem.lgv.path-step-count}
\lean{LGV1.path_stepCount_eq, LGV1.path_eastStepCount_eq}
\leanhelper
\leanok
Let $(a,b), (c,d) \in \mathbb{Z}^2$.

\textbf{Observation 1:} Each path from $(a,b)$ to $(c,d)$
has exactly $c+d-a-b$ steps.

\textbf{Observation 2:} Each path from $(a,b)$ to $(c,d)$
has exactly $c-a$ east-steps.
\end{lemma}

\begin{proof}
\textbf{Observation 1:} We define the \emph{coordinate sum} of a
point $(x,y) \in \mathbb{Z}^{2}$ to be $x+y$, denoted
$\operatorname{cs}(x,y)$. The coordinate sum increases by exactly $1$
along each arc. If the path visits vertices $v_0, v_1, \ldots, v_n$ with
$v_0 = (a,b)$ and $v_n = (c,d)$, then by the telescope principle,
\[
  n = \sum_{i=1}^n \bigl(\operatorname{cs}(v_i) - \operatorname{cs}(v_{i-1})\bigr)
    = \operatorname{cs}(c,d) - \operatorname{cs}(a,b) = c + d - a - b.
\]

\textbf{Observation 2:} The x-coordinate of a point increases by
exactly $1$ along each east-step and stays unchanged along each
north-step. By telescoping on the x-coordinate:
\[
  (\text{\# of east-steps}) = \operatorname{x}(c,d) - \operatorname{x}(a,b) = c - a.
\]
\end{proof}

\begin{proposition}
\label{prop.lgv.1-paths.ct}
\lean{LGV1.numPaths_eq_card, LGV1.numPaths, LGV1.no_paths_when_sum_decreases}
\leantarget
\leanok
Let $\left(  a,b\right)  \in\mathbb{Z}^{2}$ and
$\left(  c,d\right)  \in\mathbb{Z}^{2}$ be two points. Then,
\[
\left(  \text{\# of paths from }\left(  a,b\right)  \text{ to }\left(
c,d\right)  \right)  =%
\begin{cases}
\dbinom{c+d-a-b}{c-a}, & \text{if }c+d\geq a+b;\\
0, & \text{if }c+d<a+b.
\end{cases}
\]
\end{proposition}

\begin{proof}
Observation~1 immediately shows that no path from $(a,b)$ to $(c,d)$
exists when $c+d < a+b$ (it would need a negative number of steps),
so the count is~$0$.

Otherwise, $c + d \ge a + b$, so $c+d-a-b \in \mathbb{N}$.
Observations 1 and 2 show that every path from $(a,b)$ to $(c,d)$
has exactly $c + d - a - b$ steps, of which exactly $c - a$ are
east-steps.

\textbf{Observation 3:} Let $p$ be a path that starts at the point $(a,b)$
and has exactly $c+d-a-b$ steps. Assume that exactly $c-a$ of these
steps are east-steps. Then, this path $p$ ends at $(c,d)$.

(This follows by applying Observations 1 and 2 to the endpoint of $p$
and comparing.)

Combining Observations 1, 2 and 3, the paths from $(a,b)$ to $(c,d)$
are precisely the paths that start at $(a,b)$ and have exactly
$c+d-a-b$ steps and exactly $c-a$ east-steps. Such a path is
uniquely determined by which $c-a$ of its $c+d-a-b$ steps are
east-steps. By the bijection principle,
\[
  \left(  \text{\# of paths from }\left(  a,b\right)  \text{ to }\left(
  c,d\right)  \right)
  = \dbinom{c+d-a-b}{c-a}.
\]
\end{proof}

\subsection{Path tuples, nipats and ipats}

\begin{definition}
\label{def.lgv.path-tups}
\lean{LGV1.kVertex, LGV1.kVertex.permute, LGV1.PathTuple, LGV1.PathTuple.isNonIntersecting, LGV1.PathTuple.isIntersecting}
\leantarget
\leanok
Let $k\in\mathbb{N}$.

\textbf{(a)} A $k$\emph{-vertex} means a $k$-tuple of lattice points.

\textbf{(b)} If $\mathbf{A}=\left(  A_{1},A_{2},\ldots,A_{k}\right)  $ is a
$k$-vertex, and if $\sigma\in S_{k}$ is a permutation, then $\sigma\left(
\mathbf{A}\right)  $ shall denote the $k$-vertex $\left(  A_{\sigma\left(
1\right)  },A_{\sigma\left(  2\right)  },\ldots,A_{\sigma\left(  k\right)
}\right)  $.

\textbf{(c)} If $\mathbf{A}=\left(  A_{1},A_{2},\ldots,A_{k}\right)  $ and
$\mathbf{B}=\left(  B_{1},B_{2},\ldots,B_{k}\right)  $ are two $k$-vertices,
then a \emph{path tuple} from $\mathbf{A}$ to $\mathbf{B}$ means a $k$-tuple
$\left(  p_{1},p_{2},\ldots,p_{k}\right)  $, where each $p_{i}$ is a path from
$A_{i}$ to $B_{i}$.

\textbf{(d)} A path tuple $\left(  p_{1},p_{2},\ldots,p_{k}\right)  $ is said
to be \emph{non-intersecting} if no two of the paths $p_{1},p_{2},\ldots
,p_{k}$ have any vertex in common.
We abbreviate ``non-intersecting path tuple'' as \emph{nipat}.

\textbf{(e)} A path tuple $\left(  p_{1},p_{2},\ldots,p_{k}\right)  $ is said
to be \emph{intersecting} if it is not non-intersecting (i.e., if two of its
paths have a vertex in common).
We abbreviate ``intersecting path tuple'' as \emph{ipat}.
\end{definition}

\subsection{Helpers for path tuple finiteness}

\begin{lemma}
\label{lem.lgv.pathtup-finite}
\lean{LGV1.pathTuplesFromTo_finite}
\leanhelper
\leanok
For any two $k$-vertices $\mathbf{A}$, $\mathbf{B}$, the set of path
tuples from $\mathbf{A}$ to $\mathbf{B}$ is finite.
\end{lemma}

\begin{proof}
\leanok
The set of path tuples embeds injectively into the product
$\prod_{i=1}^k \{\text{paths from } A_i \text{ to } B_i\}$, which is finite
by Lemma~\ref{lem.lgv.paths-finite}.
\end{proof}

\begin{lemma}
\label{lem.lgv.nipats-finite}
\lean{LGV1.nipatsFromTo_finite}
\leanhelper
\leanok
For any two $k$-vertices $\mathbf{A}$, $\mathbf{B}$, the set of
nipats from $\mathbf{A}$ to $\mathbf{B}$ is finite.
\end{lemma}

\begin{proof}
\leanok
Subset of the finite set of all path tuples (Lemma~\ref{lem.lgv.pathtup-finite}).
\end{proof}

\begin{lemma}
\label{lem.lgv.ipats-finite}
\lean{LGV1.ipatsFromTo_finite}
\leanhelper
\leanok
For any two $k$-vertices $\mathbf{A}$, $\mathbf{B}$, the set of
ipats from $\mathbf{A}$ to $\mathbf{B}$ is finite.
\end{lemma}

\begin{proof}
\leanok
Subset of the finite set of all path tuples (Lemma~\ref{lem.lgv.pathtup-finite}).
\end{proof}

\subsection{Crowded vertices and the minimum crowded path index}

\begin{definition}
\label{def.lgv.crowded}
\lean{LGV1.PathTuple.isCrowded, LGV1.PathTuple.crowdedPathIndices, LGV1.PathTuple.minCrowdedPathIndex, LGV1.PathTuple.crowdedVerticesOnPath}
\leanhelper
\leanok
Let $\mathbf{p} = (p_1, \ldots, p_k)$ be a path tuple.

\textbf{(a)} A vertex $v$ is \emph{crowded} in $\mathbf{p}$ if $v$ lies on at least
two distinct paths $p_i, p_j$ with $i \ne j$.

\textbf{(b)} The \emph{crowded path indices} of $\mathbf{p}$ is the set of indices
$i$ such that $p_i$ contains a vertex shared with some other path $p_j$.

\textbf{(c)} For an ipat $\mathbf{p}$, the \emph{minimum crowded path index} is
the smallest $i$ in the crowded path indices.

\textbf{(d)} The \emph{crowded vertices on path $i$} is the set of vertices on
$p_i$ that also appear on some other path.
\end{definition}

\begin{lemma}
\label{lem.lgv.ipat-iff-crowded}
\lean{LGV1.PathTuple.isIntersecting_iff_exists_crowded, LGV1.PathTuple.isIntersecting_iff_crowdedPathIndices_nonempty}
\leanhelper
\leanok
A path tuple is intersecting if and only if it has at least one crowded vertex.
Equivalently, a path tuple is intersecting if and only if its set of crowded
path indices is nonempty.
\end{lemma}

\begin{proof}
\leanok
Immediate from the definitions: the negation of ``all pairs of paths are
disjoint'' is equivalent to the existence of a shared vertex.
\end{proof}

\subsection{The LGV lemma for two paths}

\subsubsection{Signed path tuple infrastructure}

\begin{definition}
\label{def.lgv.signed-pathtup2}
\lean{LGV1.SignedPathTuple2, LGV1.SignedPathTuple2.sign, LGV1.SignedPathTuple2.isIntersecting, LGV1.pathMatrix2, LGV1.numNipats2}
\leanhelper
\leanok
Given four lattice points $A, A', B, B'$:

\textbf{(a)} A \emph{signed path tuple} is a pair of paths $(p_0, p_1)$ together
with a sign flag indicating whether the destinations
are $(B, B')$ (sign $+1$) or $(B', B)$ (sign $-1$).

\textbf{(b)} The \emph{path count matrix} is
the $2 \times 2$ integer matrix with entries $\#\mathrm{paths}(A_i \to B_j)$.

\textbf{(c)} The number of nipats from
$(A, A')$ to $(B, B')$.
\end{definition}

\begin{lemma}
\label{lem.lgv.signed-pathtup2-finite}
\lean{LGV1.signedPathTuples2_finite, LGV1.signedNipats2_finite, LGV1.signedIpats2_finite}
\leanhelper
\leanok
The sets of signed path tuples (all, non-intersecting, and intersecting) are finite.
\end{lemma}

\begin{proof}
\leanok
The signed path tuples inject into the product of two finite path sets times
a two-element set. All three subsets are finite as subsets of this finite set.
\end{proof}

\subsubsection{First intersection point}

\begin{lemma}
\label{lem.lgv.first-intersection-idx}
\lean{LGV1.firstIntersectionIdx, LGV1.firstIntersectionIdx_lt}
\leanhelper
\leanok
Given an intersecting signed path tuple $(p_0, p_1)$, the \emph{first intersection
index} is the smallest index $i$ such that the $i$-th vertex of $p_0$ belongs
to the vertices of $p_1$. This index is strictly less than the number of
vertices of $p_0$.
\end{lemma}

\begin{proof}
\leanok
The existence of a shared vertex guarantees that a valid index
is found before the end of the vertex list.
\end{proof}

\begin{lemma}
\label{lem.lgv.first-intersection-props}
\lean{LGV1.firstIntersection, LGV1.firstIntersection_mem_path0, LGV1.firstIntersection_mem_path1, LGV1.firstIntersection_is_first}
\leanhelper
\leanok
The \emph{first intersection point} $v$ of an intersecting signed path tuple $(p_0, p_1)$
satisfies:
\begin{enumerate}
\item $v$ lies on $p_0$ (at the first intersection index).
\item $v$ lies on $p_1$.
\item $v$ is \emph{first}: no vertex of $p_0$ at an earlier index belongs to $p_1$.
\end{enumerate}
\end{lemma}

\begin{proof}
Property (1) is by definition. Property (2) follows from the search
semantics: the predicate tested is membership in $p_1$'s vertices.
Property (3) follows from the minimality of the search: all earlier
indices fail the predicate.
\end{proof}

\begin{lemma}
\label{lem.lgv.splitpath-head}
\lean{LGV1.splitPathAt}
\leanhelper
\leanok
Splitting a path at a vertex $v$ (by taking the first steps up to the
index of $v$ and the remaining steps after) produces a head ending at $v$
and a tail starting at $v$.
\end{lemma}

\begin{proof}
\leanok
By Lemma~\ref{lem.lgv.path-take-drop}, taking $n$ steps gives a path from
$A$ to the intermediate point, and dropping $n$ steps gives a path from
that point to $B$. The key fact is that the vertex at the
split position equals $v$, so the intermediate point is $v$.
\end{proof}

\begin{lemma}
\label{lem.lgv.first-intersection-preserved}
\lean{LGV1.firstIntersection_preserved}
\leanhelper
\leanok
After applying the ipat involution (exchanging tails at the first
intersection point $v$), the first intersection point of the
resulting path tuple is still $v$.

This is the key technical lemma for proving the involution is
involutive. It relies on:
\begin{itemize}
\item The prefix of $p_0$ up to the split index is unchanged.
\item $v$ is at the same index in the new path.
\item No earlier vertex of $p_0$ lies in $p_1$'s vertices (by minimality).
\item Vertices at distinct positions on a lattice path are distinct
  (because the coordinate sum strictly increases).
\end{itemize}
\end{lemma}

\begin{proof}
The new path $q_0 = \mathrm{head}(p_0) \cdot \mathrm{tail}(p_1)$ shares
the same prefix as $p_0$ up to the split index. The value
$v$ at this index is preserved.
For $j$ less than the split index, the $j$-th vertex of $q_0$
equals the $j$-th vertex of $p_0$, which is not among the new $p_1$'s
vertices (the head of old $p_1$ is a subset of old $p_1$'s
vertices, and the tail of old $p_0$ cannot contain an earlier
vertex because coordinate sums strictly increase). Thus
the search on $q_0$ returns the same index.
\end{proof}

\begin{lemma}
\label{lem.lgv.2paths.signedipat}
\lean{LGV1.ipatInvolution, LGV1.ipatInvolution_sign, LGV1.ipatInvolution_isIntersecting, LGV1.ipatInvolution_involutive, LGV1.firstIntersection, LGV1.splitPathAt}
\leanhelper
\leanok
The sign-reversing involution on intersecting path tuples for $k = 2$.

Given an intersecting path tuple $(p, p')$ in the set
$\mathcal{X}$ of ipats from $(A, A')$ to $(B, B')$ or to $(B', B)$:

\begin{enumerate}
\item Since $(p, p') \in \mathcal{X}$, the paths $p$ and $p'$ have a
  vertex in common. Let $v$ be the first one.
  (The first one on $p$ or the first one on $p'$? These are the same, since
  our digraph is acyclic.)
  We call this vertex $v$ the \emph{first intersection} of $(p, p')$.
\item Split each path at $v$ into a head (before $v$) and a tail (after $v$).
\item Exchange the tails:
  $q = (\text{head of } p) \cup (\text{tail of } p')$ and
  $q' = (\text{head of } p') \cup (\text{tail of } p)$.
\item Define $f(p, p') = (q, q')$.
\end{enumerate}

Then $f$ is a sign-reversing involution on $\mathcal{X}$ with no fixed
points.
\end{lemma}

\begin{proof}
If $(p, p')$ is an ipat from $(A, A')$ to $(B, B')$, then
$f(p, p') = (q, q')$ is an ipat from $(A, A')$ to $(B', B)$
(exchanging tails swaps the endpoints), and vice versa.
Thus $f$ is well-defined and sign-reversing.

To show $f \circ f = \operatorname{id}$: after exchanging tails, the
heads are unchanged, so $v$ is still the first intersection point
(Lemma~\ref{lem.lgv.first-intersection-preserved}).
Exchanging tails again undoes the swap.

Since $f$ is sign-reversing, it has no fixed points.
\end{proof}

\begin{lemma}
\label{lem.lgv.2paths.ipatcancel}
\lean{LGV1.sum_signedIpats2_eq_zero}
\leanhelper
\leanok
For any four lattice points $A, A', B, B'$,
\[
  \sum_{(p,p') \in \mathcal{X}} \operatorname{sign}(p, p') = 0,
\]
where $\mathcal{X}$ is the set of intersecting path tuples and the sign
is $+1$ for path tuples to $(B, B')$ and $-1$ for path tuples to
$(B', B)$.
\end{lemma}

\begin{proof}
By Lemma~\ref{lem.sign.cancel2} applied to the sign-reversing
involution of Lemma~\ref{lem.lgv.2paths.signedipat}.
\end{proof}

\begin{lemma}
\label{lem.lgv.2paths.partition}
\lean{LGV1.sum_signedPathTuples2, LGV1.sum_signedNipats2}
\leanhelper
\leanok
The sum over all signed path tuples decomposes as:
\begin{align*}
  &\bigl(\text{\# of path tuples from } (A,A') \text{ to }
   (B,B')\bigr)
   - \bigl(\text{\# of path tuples from } (A,A') \text{ to }
   (B',B)\bigr) \\
  &= \sum_{(p,p') \in \mathcal{A}} \operatorname{sign}(p,p'),
\end{align*}
and the sum over nipats gives:
\begin{align*}
  \sum_{(p,p') \in \mathcal{A} \setminus \mathcal{X}}
    \operatorname{sign}(p,p')
  &= \bigl(\text{\# of nipats from } (A,A') \text{ to } (B,B')\bigr) \\
  &\quad - \bigl(\text{\# of nipats from } (A,A') \text{ to } (B',B)\bigr).
\end{align*}
\end{lemma}

\begin{proof}
The first identity follows from the product rule: the product of path
counts equals the count of path tuples.

The second identity is immediate from the definitions: among the nipats,
those going to $(B,B')$ contribute $+1$ each and those going to
$(B',B)$ contribute $-1$ each.
\end{proof}

\begin{proposition}
[LGV lemma for two paths]\label{prop.lgv.2paths.count}
\lean{LGV1.lgv_two_paths}
\leantarget
\leanok
Let $\left(
A,A^{\prime}\right)  $ and $\left(  B,B^{\prime}\right)  $ be two $2$-vertices
(i.e., let $A,A^{\prime},B,B^{\prime}$ be four lattice points). Then,
\begin{align*}
&  \det\left(
\begin{array}
[c]{cc}%
\left(  \text{\# of paths from }A\text{ to }B\right)  & \left(  \text{\# of
paths from }A\text{ to }B^{\prime}\right) \\
\left(  \text{\# of paths from }A^{\prime}\text{ to }B\right)  & \left(
\text{\# of paths from }A^{\prime}\text{ to }B^{\prime}\right)
\end{array}
\right) \\
&  =\left(  \text{\# of nipats from }\left(  A,A^{\prime}\right)  \text{ to
}\left(  B,B^{\prime}\right)  \right) \\
&  \ \ \ \ \ \ \ \ \ \ -\left(  \text{\# of nipats from }\left(  A,A^{\prime
}\right)  \text{ to }\left(  B^{\prime},B\right)  \right)  .
\end{align*}
\end{proposition}

\begin{proof}
We have
\begin{align*}
&  \det\left(
\begin{array}
[c]{cc}%
\left(  \text{\# of paths from }A\text{ to }B\right)  & \left(  \text{\# of
paths from }A\text{ to }B^{\prime}\right) \\
\left(  \text{\# of paths from }A^{\prime}\text{ to }B\right)  & \left(
\text{\# of paths from }A^{\prime}\text{ to }B^{\prime}\right)
\end{array}
\right) \\
&  =\left(  \text{\# of path tuples from }\left(  A,A^{\prime}\right)  \text{
to }\left(  B,B^{\prime}\right)  \right) \\
&  \ \ \ \ \ \ \ \ \ \ -\left(  \text{\# of path tuples from }\left(
A,A^{\prime}\right)  \text{ to }\left(  B^{\prime},B\right)  \right)
\end{align*}
(by the product rule). By Lemma~\ref{lem.lgv.2paths.partition},
the right hand side equals
$\sum_{(p,p') \in \mathcal{A}} \operatorname{sign}(p,p')$.

By Lemma~\ref{lem.sign.cancel2}, using the sign-reversing involution
of Lemma~\ref{lem.lgv.2paths.signedipat} (which uses
Lemma~\ref{lem.lgv.2paths.ipatcancel}), we obtain
\[
\sum_{\left(  p,p^{\prime}\right)  \in\mathcal{A}}\operatorname{sign}%
\left(  p,p^{\prime}\right)  =\sum_{\left(  p,p^{\prime}\right)  \in\mathcal{A}%
\setminus\mathcal{X}}\operatorname{sign}\left(  p,p^{\prime}\right).
\]
The right hand side equals
\begin{align*}
&  \left(  \text{\# of nipats from }\left(  A,A^{\prime}\right)  \text{ to
}\left(  B,B^{\prime}\right)  \right) \\
&  \ \ \ \ \ \ \ \ \ \ -\left(  \text{\# of nipats from }\left(  A,A^{\prime
}\right)  \text{ to }\left(  B^{\prime},B\right)  \right)
\end{align*}
by Lemma~\ref{lem.lgv.2paths.partition}.
\end{proof}

\subsection{The baby Jordan curve theorem}

\begin{proposition}
[baby Jordan curve theorem]\label{prop.lgv.jordan-2}
\lean{LGV1.baby_jordan_curve, LGV1.no_nipats_under_nw, LGV1.isWeaklyNorthwestOf}
\leantarget
\leanok
Let $A$, $B$, $A^{\prime}$
and $B^{\prime}$ be four lattice points satisfying
\begin{align*}
\operatorname{x}\left(  A^{\prime}\right)   &  \leq\operatorname{x}\left(
A\right)  ,\ \ \ \ \ \ \ \ \ \ \operatorname{y}\left(  A^{\prime}\right)
\geq\operatorname{y}\left(  A\right)  ,\\
\operatorname{x}\left(  B^{\prime}\right)   &  \leq\operatorname{x}\left(
B\right)  ,\ \ \ \ \ \ \ \ \ \ \operatorname{y}\left(  B^{\prime}\right)
\geq\operatorname{y}\left(  B\right)  .
\end{align*}
(That is, $A'$ is weakly northwest of $A$, and $B'$ is weakly northwest
of $B$.)

Let $p$ be any path from $A$ to $B^{\prime}$. Let $p^{\prime}$ be any path
from $A^{\prime}$ to $B$. Then, $p$ and $p^{\prime}$ have a vertex in common.

In particular, there are no nipats from $(A, A')$ to $(B', B)$.
\end{proposition}

\begin{proof}
The proof is by strong induction on $\ell(p) + \ell(p')$, the sum of
path lengths.

\textbf{Case 3} ($A' = A$): Then $A$ is a common vertex of both
paths and we are done.

\textbf{Case 1} ($\operatorname{y}(A') > \operatorname{y}(A)$):
Let $P$ be the second vertex of $p$. Every vertex of the
subpath $r$ from $P$ to $B'$ is a vertex of $p$, so $r$ and $p'$
have no vertex in common. Also $\ell(r) + \ell(p') < \ell(p) + \ell(p')$.
If $\operatorname{y}(A') \ge \operatorname{y}(P)$, then we could apply
the induction hypothesis to $r$ and $p'$, contradicting $r \cap p' = \emptyset$.
Hence $\operatorname{y}(A') < \operatorname{y}(P)$, so
$\operatorname{y}(A') \le \operatorname{y}(P) - 1$. If the arc from $A$ to $P$
were an east-step, we would have $\operatorname{y}(P) = \operatorname{y}(A)$,
contradicting $\operatorname{y}(A) \le \operatorname{y}(A') < \operatorname{y}(P)$.
So the arc is a north-step, giving $\operatorname{y}(P) = \operatorname{y}(A) + 1$,
hence $\operatorname{y}(A') \le \operatorname{y}(A)$,
contradicting $\operatorname{y}(A') > \operatorname{y}(A)$.

\textbf{Case 2} ($\operatorname{x}(A') < \operatorname{x}(A)$): Symmetric
to Case 1, looking at the first step of $p'$ and using reflection
across the line $x = y$.

The consequence that there are no nipats from $(A, A')$ to $(B', B)$
follows immediately: if any pair of paths (one from $A$ to $B'$, one from
$A'$ to $B$) must intersect, then there can be no nipats from $(A, A')$
to $(B', B)$.
\end{proof}

\subsection{Log-concavity of binomial coefficients}

\begin{corollary}
\label{cor.lgv.binom-unimod}
\lean{LGV1.binom_log_concave}
\leantarget
\leanok
Let $n,k\in\mathbb{N}$. Then, $\dbinom{n}{k}%
^{2}\geq\dbinom{n}{k-1}\cdot\dbinom{n}{k+1}$.
\end{corollary}

\begin{proof}
\leanok
\emph{Combinatorial proof (from the source text, using LGV):}
Define four lattice points
$A=\left(  1,0\right)  $ and $A^{\prime}=\left(  0,1\right)  $ and $B=\left(
k+1,n-k\right)  $ and $B^{\prime}=\left(  k,n-k+1\right)  $. Then, Proposition
\ref{prop.lgv.2paths.count} yields
\begin{align*}
&  \det\left(
\begin{array}
[c]{cc}%
\left(  \text{\# of paths from }A\text{ to }B\right)  & \left(  \text{\# of
paths from }A\text{ to }B^{\prime}\right) \\
\left(  \text{\# of paths from }A^{\prime}\text{ to }B\right)  & \left(
\text{\# of paths from }A^{\prime}\text{ to }B^{\prime}\right)
\end{array}
\right) \\
&  =\left(  \text{\# of nipats from }\left(  A,A^{\prime}\right)  \text{ to
}\left(  B,B^{\prime}\right)  \right)  -\underbrace{\left(  \text{\# of nipats
from }\left(  A,A^{\prime}\right)  \text{ to }\left(  B^{\prime},B\right)
\right)  }_{=0}\\
&  =\left(  \text{\# of nipats from }\left(  A,A^{\prime}\right)  \text{ to
}\left(  B,B^{\prime}\right)  \right)  \geq0.
\end{align*}
(Here, the second term is $0$ by Proposition \ref{prop.lgv.jordan-2}.)
However, Proposition \ref{prop.lgv.1-paths.ct} yields
\[
\left(
\begin{array}
[c]{cc}%
\left(  \text{\# of paths from }A\text{ to }B\right)  & \left(  \text{\# of
paths from }A\text{ to }B^{\prime}\right) \\
\left(  \text{\# of paths from }A^{\prime}\text{ to }B\right)  & \left(
\text{\# of paths from }A^{\prime}\text{ to }B^{\prime}\right)
\end{array}
\right)
=\left(
\begin{array}
[c]{cc}%
\dbinom{n}{k} & \dbinom{n}{k-1}\\
\dbinom{n}{k+1} & \dbinom{n}{k}%
\end{array}
\right),
\]
so the determinant is $\dbinom{n}{k}^{2}-\dbinom{n}%
{k-1}\cdot\dbinom{n}{k+1}\geq 0$.

\emph{Lean formalization (algebraic proof):} The Lean proof avoids the
LGV machinery and proceeds purely algebraically, using the recurrences
\begin{align*}
  \binom{n}{k+1}(k+1) &= \binom{n}{k}(n-k), \\
  \binom{n}{k} \cdot k &= \binom{n}{k-1}(n-k+1).
\end{align*}
These allow rewriting both sides after multiplying by $k(k+1) > 0$.
The inequality reduces to $(n-k+1)(k+1) \ge (n-k) \cdot k$, which
holds because the difference is $n + 1 \ge 0$. This proof is
complete.
\end{proof}

\subsection{The LGV lemma for $k$ paths}

\subsubsection{Definitions for general $k$}

\begin{definition}
\label{def.lgv.pathmatrixK}
\lean{LGV1.pathMatrixK, LGV1.numNipatsK, LGV1.numIpatsK}
\leanhelper
\leanok
Let $k \in \mathbb{N}$ and let $\mathbf{A}, \mathbf{B}$ be two $k$-vertices.

\textbf{(a)} The \emph{path count matrix} $M$
is the $k \times k$ integer matrix with entries
$M_{i,j} = \#\mathrm{paths}(A_i \to B_j)$.

\textbf{(b)} For $\sigma \in S_k$, the number of nipats from $\mathbf{A}$ to $\sigma(\mathbf{B})$.

\textbf{(c)} For $\sigma \in S_k$, the number of ipats from $\mathbf{A}$ to $\sigma(\mathbf{B})$.
\end{definition}

\subsubsection{Bridge between LGV1 and LGV2}

\begin{lemma}
\label{lem.lgv.step-equiv}
\lean{LGV1.latticeStepEquiv, LGV1.latticePathToLatticePath', LGV1.latticePath'ToLatticePath, LGV1.latticePathToLatticePath'_endpoint}
\leanhelper
\leanok
The lattice step types in LGV1 and LGV2 are equivalent:
the two representations of lattice steps are isomorphic, and this
equivalence preserves the step application function. Lattice paths can be
converted between representations, preserving endpoints and vertices.
\end{lemma}

\begin{proof}
\leanok
Both types are two-element enumerations with the same constructors.
The equivalence and endpoint preservation follow by case analysis.
\end{proof}

\begin{lemma}
\label{lem.lgv.numipats-eq-ipatfinset}
\lean{LGV1.numIpatsK_eq_ipatFinset_card}
\leanhelper
\leanok
The ipat count (defined via one representation of path tuples)
equals the cardinality of the ipat set (defined via the other representation).
This is the key bridge that allows the proof to invoke the sign-reversing
involution from the weighted LGV formalization.
\end{lemma}

\begin{proof}
The proof constructs an equivalence between the two representations of
intersecting path tuples, converting between list-of-steps and
list-of-vertices representations,
then verifies that the intersection property is preserved at each step.
\end{proof}

\subsubsection{Product rule and partition lemmas}

\begin{lemma}
\label{lem.lgv.kpaths.partition}
Let $k \in \mathbb{N}$ and let $\mathbf{A}, \mathbf{B}$ be two
$k$-vertices. Then:

\textbf{(a)} (Product rule) For any $\sigma \in S_k$,
\[
  \prod_{i=1}^k \bigl(\text{\# of paths from } A_i \text{ to }
  B_{\sigma(i)}\bigr) = \bigl(\text{\# of path tuples from }
  \mathbf{A} \text{ to } \sigma(\mathbf{B})\bigr).
\]

\textbf{(b)} (Partition) Each path tuple is either a nipat or an ipat,
so
\[
  \bigl(\text{\# of path tuples}\bigr) = \bigl(\text{\# of nipats}\bigr)
  + \bigl(\text{\# of ipats}\bigr).
\]

\textbf{(c)} (Determinant expansion)
\[
  \det\bigl(\bigl(\text{\# of paths from } A_i \text{ to }
  B_j\bigr)_{1 \le i,j \le k}\bigr) = \sum_{\sigma \in S_k} (-1)^\sigma
  \prod_{i=1}^k \bigl(\text{\# of paths from } A_i \text{ to }
  B_{\sigma(i)}\bigr).
\]
\end{lemma}

\begin{proof}
\textbf{(a)} A path tuple from $\mathbf{A}$ to $\sigma(\mathbf{B})$ is
just a tuple of independent choices (one path per index), so by the
product rule the count is the product.

\textbf{(b)} Every path tuple is intersecting or not (classical logic);
the sets are disjoint and their union is the set of all path tuples.

\textbf{(c)} This is the Leibniz formula
$\det(M) = \sum_\sigma (-1)^\sigma \prod_i M_{i,\sigma(i)}$.
\end{proof}

\begin{lemma}
\label{lem.lgv.kpaths.involution}
The ipat cancellation lemma for general $k$:
\[
  \sum_{\sigma \in S_k} (-1)^\sigma \cdot
  \bigl(\text{\# of ipats from } \mathbf{A} \text{ to }
  \sigma(\mathbf{B})\bigr) = 0.
\]
This is proved by a sign-reversing involution on pairs
$(\sigma, \mathbf{p})$ where $\mathbf{p}$ is an ipat from $\mathbf{A}$
to $\sigma(\mathbf{B})$. The involution:
\begin{enumerate}
\item A point $u$ is \emph{crowded} if it is contained in at least
  two of the paths $p_1, p_2, \ldots, p_k$.
\item Pick the smallest $i \in [k]$ such that $p_i$ contains a crowded
  point.
\item Let $v$ be the first crowded point on $p_i$.
\item Let $j > i$ be the largest index such that $v \in p_j$.
\item Exchange the tails of $p_i$ and $p_j$ at $v$.
\item Replace $\sigma$ by $\sigma \cdot t_{i,j}$ (composing with the
  transposition swapping $i$ and $j$), which flips the sign
  (by Proposition~\ref{prop.perm.sign.props}~\textbf{(b)} and
  \textbf{(d)}).
\end{enumerate}
The map $f : (\sigma, \mathbf{p}) \mapsto (\sigma', \mathbf{q})$ is
an involution because the heads of $p_i$ and $p_j$ are unchanged, so
the same $v$, $i$, $j$ are selected again.
\end{lemma}

\begin{proof}
The proof uses:
\begin{itemize}
\item The partition of path tuples (Lemma~\ref{lem.lgv.kpaths.partition}).
\item The sign-reversing involution described above.
\item Lemma~\ref{lem.sign.cancel2} to cancel the signed sum of ipats.
\item The bridge Lemma~\ref{lem.lgv.numipats-eq-ipatfinset} to invoke the
  sign-reversing involution from the weighted LGV formalization.
\end{itemize}
\end{proof}

\begin{proposition}
[LGV lemma, lattice counting version]\label{prop.lgv.kpaths.count}
\lean{LGV1.lgv_k_paths}
\leantarget
\leanok
Let
$k\in\mathbb{N}$. Let $\mathbf{A}=\left(  A_{1},A_{2},\ldots,A_{k}\right)  $
and $\mathbf{B}=\left(  B_{1},B_{2},\ldots,B_{k}\right)  $ be two
$k$-vertices. Then,
\begin{align*}
&  \det\left(  \left(  \text{\# of paths from }A_{i}\text{ to }B_{j}\right)
_{1\leq i\leq k,\ 1\leq j\leq k}\right) \\
&  =\sum_{\sigma\in S_{k}}\left(  -1\right)  ^{\sigma}\left(  \text{\# of
nipats from }\mathbf{A}\text{ to }\sigma\left(  \mathbf{B}\right)  \right)  .
\end{align*}
\end{proposition}

\begin{proof}
Define
$\mathcal{A}:=\left\{  \left(  \sigma,\mathbf{p}\right)  \ \mid\ \sigma\in
S_{k}\text{, and}\ \mathbf{p}\text{ is a path tuple from }\mathbf{A}\text{ to
}\sigma\left(  \mathbf{B}\right)  \right\}$
and
$\mathcal{X}:=\left\{  \left(
\sigma,\mathbf{p}\right)  \in\mathcal{A}\ \mid\ \mathbf{p}\text{ is
intersecting}\right\}$.
Set
$\operatorname{sign}\left(  \sigma,\mathbf{p}\right)  :=\left(  -1\right)
^{\sigma}$.

By Lemma~\ref{lem.sign.cancel2} applied to the sign-reversing involution
of Lemma~\ref{lem.lgv.kpaths.involution},
\[
\sum_{\left(  \sigma,\mathbf{p}\right)  \in\mathcal{A}}\left(  -1\right)
^{\sigma}=\sum_{\left(  \sigma,\mathbf{p}\right)  \in\mathcal{A}%
\setminus\mathcal{X}}\left(  -1\right)  ^{\sigma}.
\]
The left hand side equals
\begin{align*}
\sum_{\sigma\in S_{k}}\left(  -1\right)  ^{\sigma}\prod_{i=1}^{k}\left(
\text{\# of paths from }A_{i}\text{ to }B_{\sigma\left(  i\right)  }\right)
&  =\det\left(  \left(  \text{\# of paths from }A_{i}\text{ to }B_{j}\right)
_{1\leq i\leq k,\ 1\leq j\leq k}\right)
\end{align*}
(by Lemma~\ref{lem.lgv.kpaths.partition}(a) and the definition of the
determinant), and the right hand side equals
\[
\sum_{\sigma\in S_{k}}\left(  -1\right)  ^{\sigma}\left(  \text{\# of
nipats from }\mathbf{A}\text{ to }\sigma\left(  \mathbf{B}\right)  \right).
\]
\end{proof}
